\section{Appendix}

\subsection{Authorship and Contributions}

This report is partitioned so each major part has a primary author responsible for writing and technical correctness.

\subsubsection*{Deniz Isikli (s215818)}
Primary author of:
\begin{itemize}
    \item Chapter 3 (Design Choices and Information-Flow Argument),
    \item Chapter 5 (Security Analysis of API Commands),
\end{itemize}

\subsubsection*{Benas Skripkiunas (s253140)}
Primary author of:
\begin{itemize}
    \item Chapter 2 (Participants, Roles, and Assumptions),
    \item Chapter 6 (Declassification Justification),
\end{itemize}

\subsubsection*{Abid Hasan (s252729)}
Primary author of:
\begin{itemize}
    \item Chapter 1 (Introduction),
    \item Chapter 4 (Server API and Security Semantics),
\end{itemize}

\subsection{Appendix A: Flow Checker Pseudocode}

The server's type checker implements the following core operations, referenced in Chapters 3, 4, and 5.

\subsubsection*{Core Flow Checking Functions}

The type checker implements: \texttt{canFlow}($L_1, L_2$) checks if $L_1 \sqsubseteq L_2$; \texttt{join}($L_1, L_2$) computes the least upper bound by taking reader-set intersections for confidentiality and writer-set unions for integrity; \texttt{check\_assignment} verifies $pc \sqcup L_e \sqsubseteq L_{target}$ for each assignment. When a branch guards an assignment, $pc$ is raised to $pc \sqcup L_{guard}$ inside the branch body.

\subsection{Appendix B: Label Comparison Examples}

Below are five concrete examples of flow checking, worked through by hand:

\textbf{Example 1: Simple Confidentiality Flow}

Source: $L_1 = (doctor:\{doctor,nurse\})$
Destination: $L_2 = (doctor:\{nurse\})$

Is $L_1 \sqsubseteq L_2$?

\begin{itemize}
    \item Confidentiality: $(doctor:\{doctor,nurse\})$ in $L_1$; $(doctor:\{nurse\})$ in $L_2$. Check: $\{doctor, nurse\} \supseteq \{nurse\}$? Yes.
    \item Integrity: neither label has integrity policies.
\end{itemize}

Result: \textbf{Yes}, $L_1 \sqsubseteq L_2$. Information can flow.

\subsection{Appendix B: Worked API Trace Example}

\textbf{Scenario:} A patient uploads medical data and the hospital declassifies a summary for research.

\textbf{Step 1: UPLOAD (Patient uploads image)}

\begin{verbatim}
Input:
  token = <valid auth for patient>
  payload = <medical image bytes>
  C = (patient:{patient, doctor})
  I = (patient <- {patient})

Server checks:
  - token is valid ✓
  - C and I parse correctly ✓
  
Server effect:
  creates object:
    dataId = 1001
    label = ((patient:{patient,doctor}), (patient <- {patient}))
    owner = patient
    payload = <image>
    
Response:
  ok(1001)
\end{verbatim}

\textbf{Step 2--3: UPLOAD\_PROGRAM and RUN\_PROGRAM}\n\nThe doctor uploads an encryption program with interface input: $(patient:\{patient,doctor\})$, output: $(patient:\{patient,doctor\})$. The server validates the flow and executes the program on the patient's data.
  - program interface: input requires (patient:{patient,doctor}) ✓
  - input labels: L_1001 = ((patient:{patient,doctor}), (patient <- {patient}))
  - required output lower bound: L_req = L_1001 = ((patient:{patient,doctor}), 
    (patient <- {patient}))
  - Is L_req ⊆ ((patient:{patient,doctor}), (patient <- {patient,doctor}))?
    Conf: (patient:{patient,doctor}) ⊆ (patient:{patient,doctor})? Yes ✓
    Integ: for (patient <- {patient,doctor}) in output, must exist 
           (patient <- w) in input with w ⊇ {patient,doctor}.
           Input has (patient <- {patient}). Is {patient} ⊇ {patient,doctor}? 
           No. ✗
    
Server response:
  error(flow_violation, "output integrity too high; 
  input is not trusted by enough writers")
\end{verbatim}

The doctor corrects the output label:

\begin{verbatim}
Input (corrected):
  outI = (patient <- {patient})

Server checks:
  - Is L_req ⊆ output? Yes ✓
  
Server effect:
  executes program in sandbox
  result bytes = <encrypted image>
  creates output object:
    dataId = 1002
    label = ((patient:{patient,doctor}), (patient <- {patient}))
    
Response:
  ok(1002)
\end{verbatim}

\textbf{Step 4: DECLASSIFY\_AND\_RELEASE (Hospital releases summary)}

\begin{verbatim}
Input:
  token = <valid auth for admin>
  dataId = 1002
  transform = "aggregate"
  params = {k = 5}
  targetC = (patient:{patient,doctor,researcher}) 
            wedge (hospital:{hospital,researcher})

Server checks:
  - 1002 exists, label is ((patient:{patient,doctor}), 
    (patient <- {patient})) ✓
  - transform "aggregate" is approved ✓
  - params k=5 satisfies k_min=5? Yes ✓
  - authority check:
    * for (patient:{patient,doctor}) → (patient:{patient,doctor,researcher}),
      admin must have authority from patient. Does admin have delegation?
      (not shown, assume yes in this example) ✓
    * for (hospital:{hospital}) → (hospital:{hospital,researcher}),
      admin must have authority from hospital. (admin typically has this) ✓
  
Server effect:
  applies aggregate transform:
  result = <anonymized aggregate summary>
  creates output object:
    dataId = 1003
    label = ((patient:{patient,doctor,researcher}), 
             (hospital:{hospital,researcher}))
  writes audit entry:
    (actor: admin, time: 2024-04-22 10:30, 
     source: 1002, target: 1003, 
     transform: aggregate, k: 5,
     authority: patient-delegation-token-xyz)
     
Response:
  ok(1003)
\end{verbatim}

\textbf{Invariant Verification:}

After all four commands:
\begin{itemize}
    \item \textbf{I1 (Label Validity)}: Objects 1001, 1002, 1003 all have valid labels. ✓
    \item \textbf{I2 (Read Safety)}: Only patients and doctors can read 1001/1002. Only patients/doctors/researchers can read 1003. Read checks enforce this. ✓
    \item \textbf{I3 (Write Safety)}: Only patients can write to 1001/1002. Integrity labels are respected. ✓
    \item \textbf{I4 (Program Safety)}: Program 2001 was flow-checked before upload; RUN\_PROGRAM verified label compatibility. ✓
    \item \textbf{I5 (Declassification Safety)}: Declassification in step 4 was authority-checked and audited. ✓
\end{itemize}
